\section{Geometrías amplituhedrales de rango uno y compatibilidad de residuos}
\label{sec:k-poligono-amplituhedral}

El registro dodecafásico permite reunir explícitamente datos positivos,
dominio grassmanniano, imagen cinemática, cobertura, grado sobre las celdas
y forma canónica. El primer resultado concierne al amplituedro
bidimensional de parámetros \((n,k,m)=(13,1,2)\); después se prolonga la
construcción a los polítopos cíclicos de rango uno y dimensión superior.
Se trata de una realización geométrica posterior del registro producido
por APP--TRIT--TPK. La operación adicional declarada es la elevación por
potencias de una partición racional ordenada. Esta elección fija una
realización particular; las correspondencias físicas posteriores conservan
sus lectores y condiciones propios.

La procedencia del registro conserva la composición
\[
 x_{\mathrm{term}}\longmapsto\mathcal R_{12}(x_{\mathrm{term}})
 \longmapsto(b^{90},b^{120},Q_{\mathrm{TPK}})
 \longmapsto U_{\mathrm{sgn}}\longmapsto K.
\]
APP conserva ambas hojas y sus residuos y cocientes; TRIT orienta el
régimen de cada transición; TPK transporta y actualiza el estado con
acarreo, frontera, ruta y memoria. La representación \(K\) se lee en el mismo
sistema de estados enriquecidos que genera conjuntamente el continuo.
Su extracción firmada y su inversión
de Hadamard producen
\begin{equation}
 K=(234,543,140,729,659,824,621,58,914,794,146,601),
 \qquad S_K=\sum_{j=1}^{12}K_j=6263.
 \label{eq:polygon-K}
\end{equation}
La prueba geométrica utiliza la positividad de esas doce coordenadas, su
orden y su total. Sus coeficientes proceden de aquella extracción y no de
la selección retrospectiva de un valor físico.

\subsection{Partición racional y elevación positiva}

Definimos trece puntos racionales de partición:
\begin{equation}
 t_0=0,\qquad t_i=\frac{K_1+\cdots+K_i}{S_K}\quad(1\leq i\leq12),
 \qquad Z_i=(1,t_i,t_i^2)^\mathsf T,\qquad Z=(Z_0\ \cdots\ Z_{12}).
 \label{eq:polygon-Z}
\end{equation}
Sus numeradores, con denominador común \(6263\), son
\[
 (0,234,777,917,1646,2305,3129,3750,3808,4722,5516,5662,6263).
\]
Así, \(0=t_0<t_1<\cdots<t_{12}=1\). La partición etiquetada y el total
conservan el registro mediante la inversa exacta
\begin{equation}
 K_i=S_K(t_i-t_{i-1}).
 \label{eq:polygon-recover-K}
\end{equation}
La partición conserva las proporciones; retener \(S_K\) conserva también
la escala entera.

\begin{proposition}[Datos cinemáticos estrictamente positivos]
\label{prop:polygon-vandermonde}
La matriz \(Z\) tiene rango tres y todos sus menores ordenados de tamaño
tres son positivos. Para \(0\leq a<b<c\leq12\),
\begin{equation}
 \langle abc\rangle:=\det(Z_a,Z_b,Z_c)
 =(t_b-t_a)(t_c-t_a)(t_c-t_b)>0.
 \label{eq:polygon-minors}
\end{equation}
\end{proposition}
\begin{proof}
Se resta la primera columna de las dos últimas y se desarrolla por la
primera fila. Cada diferencia \(t_j^2-t_a^2\) factoriza como
\((t_j-t_a)(t_j+t_a)\), lo que da el producto indicado. Las sumas parciales
estrictamente crecientes prueban su positividad. Cualquiera de esos
menores acredita el rango tres.
\end{proof}

\subsection{Imagen, fibras y celdas de grado uno}

El Grassmanniano \(\operatorname{Gr}_{\geq0}(1,13)\) es el espacio
proyectivo de vectores \(c=(c_0,\ldots,c_{12})\) no negativos y no nulos,
módulo escala positiva. La coordenatización \(\sum c_j=1\) lo identifica con el
símplex cerrado de dimensión doce. Definimos
\begin{equation}
 \Phi_Z:\operatorname{Gr}_{\geq0}(1,13)\longrightarrow\mathbb P^2,
 \qquad[c]\longmapsto\left[\sum_{j=0}^{12}c_jZ_j\right].
 \label{eq:polygon-map}
\end{equation}
En la coordenatización \(Y_0=1\), sean \(v_j=(t_j,t_j^2)\) y
\(P_K=\operatorname{conv}(v_0,\ldots,v_{12})\).

\begin{theorem}[Amplituedro poligonal del registro]
\label{thm:polygon-amplituhedron}
La imagen cerrada de \(\Phi_Z\) es exactamente el polígono convexo
\(P_K=\mathcal A_{13,1,2}(Z)\), de dimensión dos y trece vértices. La imagen
del Grassmanniano estrictamente positivo es el interior de \(P_K\).
Las once celdas
\begin{equation}
 \mathcal S_i=\{[c]:c_0,c_i,c_{i+1}>0,\
 c_j=0\text{ si }j\notin\{0,i,i+1\}\},\qquad 1\leq i\leq11,
 \label{eq:polygon-cells}
\end{equation}
se aplican biyectivamente, con grado uno y orientación positiva, sobre los
interiores de \(T_i=[v_0,v_i,v_{i+1}]\). Sus clausuras cubren \(P_K\) y sus
interiores son disjuntos dos a dos.
\end{theorem}
\begin{proof}
En la normalización \(\sum c_j=1\), la aplicación es
\(c\mapsto\sum c_jv_j\), cuya imagen es la envolvente convexa. Los menores
de \eqref{eq:polygon-minors} muestran que cada triple ordenado tiene
orientación positiva. Cada lado consecutivo deja todos los otros vértices
estrictamente a su izquierda; también lo hace el lado final de \(v_{12}\)
a \(v_0\). Todos los puntos son, por tanto, vértices extremos de un polígono
convexo.

Toda combinación con \(c_j>0\) es interior, pues en cualquier recta de
soporte alguno de los vértices está estrictamente en el semiplano
interior. Recíprocamente, sean \(y\) interior y
\(b=13^{-1}\sum v_j\). Para \(\varepsilon>0\) suficientemente pequeño,
\(y'=(y-\varepsilon b)/(1-\varepsilon)\) permanece en el polígono.
Escribiendo \(y'=\sum a_jv_j\), \(a_j\geq0\), \(\sum a_j=1\), se obtiene
\(y=\sum((1-\varepsilon)a_j+\varepsilon/13)v_j\), con todos los
coeficientes positivos.

Las diagonales desde \(v_0\) dividen el polígono en los once triángulos.
Sobre cada soporte triple, las coordenadas baricéntricas son únicas,
porque \(\langle0,i,i+1\rangle>0\); su inversa explícita aparece en
\eqref{eq:polygon-barycentric}. Cada soporte define una celda positroidal
de rango uno con coordenadas \([1:a:b]\), \(a,b>0\), en esas posiciones.
Una red dirigida planar con una fuente y tres trayectorias de pesos
\(1,a,b\) hacia los tres destinos realiza esas coordenadas por medición
de frontera. El determinante positivo prueba orientación y grado uno.
\end{proof}

La fibra genérica de la aplicación completa tiene dimensión diez:
\(\sum c_j=1\), \(\sum c_jt_j=x\) y \(\sum c_jt_j^2=y\) son tres
ecuaciones independientes en trece coordenadas. Para \(y\) interior la
fibra contiene un punto estrictamente positivo, por lo que su intersección
con el símplex tiene dimensión \(13-3=10\). El grado uno concierne a las
celdas triangulares; la forma bidimensional se obtiene mediante esas
celdas.

\subsection{Forma canónica y residuos orientados}

Para \(Y=(1,x,y)^\mathsf T\), definimos
\begin{equation}
 \ell_{ab}(Y)=\det(Z_a,Z_b,Y)
 =(t_b-t_a)\{y-(t_a+t_b)x+t_at_b\}.
 \label{eq:polygon-lines}
\end{equation}
La expresión vale también cuando \(b<a\). Para \(a<b<c\), sean
\(D_{abc}=\langle abc\rangle\) y
\begin{equation}
 \lambda_a=\frac{\ell_{bc}}{D_{abc}},\qquad
 \lambda_b=\frac{\ell_{ca}}{D_{abc}},\qquad
 \lambda_c=\frac{\ell_{ab}}{D_{abc}},\qquad
 \lambda_a+\lambda_b+\lambda_c=1.
 \label{eq:polygon-barycentric}
\end{equation}
La orientación afín es \(dx\wedge dy\). Se fija la convención de frontera
\(\operatorname{Res}^{\partial}(\eta\wedge du/u)=\eta|_{u=0}\):
la diferencial normal ocupa la última posición. En dimensión dos, el
residuo de Poincaré que coloca \(du/u\) primero tiene el signo opuesto.
Las cancelaciones siguientes usan de forma uniforme la convención declarada.

\begin{theorem}[Forma racional y cancelación de las diagonales]
\label{thm:polygon-form}
La forma canónica del triángulo orientado \([v_a,v_b,v_c]\) es
\begin{equation}
 \Omega_{abc}=
 \frac{D_{abc}^2\,dx\wedge dy}
 {\ell_{ab}(Y)\ell_{bc}(Y)\ell_{ca}(Y)}.
 \label{eq:polygon-triangle-form}
\end{equation}
La forma canónica del polígono es
\begin{equation}
 \boxed{\displaystyle
 \Omega_{P_K}(x,y)=
 \sum_{i=1}^{11}
 \frac{t_it_{i+1}(t_{i+1}-t_i)\,dx\wedge dy}
 {(y-t_ix)\{y-(t_i+t_{i+1})x+t_it_{i+1}\}(t_{i+1}x-y)}.}
 \label{eq:polygon-full-form}
\end{equation}
Los polos de las diagonales interiores se cancelan. Los únicos divisores
polares son las trece rectas de los lados exteriores. Parametrizado cada
lado desde su vértice inicial al final por \(s\in[0,1]\), su residuo de
frontera es \(ds/[s(1-s)]\).
\end{theorem}
\begin{proof}
La forma del símplex baricéntrico es
\(d\lambda_b\wedge d\lambda_c/(\lambda_a\lambda_b\lambda_c)\).
El jacobiano satisface
\(dx\wedge dy=D_{abc}\,d\lambda_b\wedge d\lambda_c\).
Sustituir \eqref{eq:polygon-barycentric} da
\eqref{eq:polygon-triangle-form}. Para \(a=0\), cancelar los tres factores
escalares de las rectas da cada sumando de \eqref{eq:polygon-full-form}.

En el lado \(ab\), pongamos \(\lambda_c=u\),
\(\lambda_b=(1-u)s\), \(\lambda_a=(1-u)(1-s)\). Entonces
\[
 \Omega_{abc}=\frac{ds}{s(1-s)}\wedge\frac{du}{u(1-u)}.
\]
El residuo es la forma del intervalo orientado. Dos triángulos contiguos
inducen orientaciones contrarias sobre su diagonal y producen residuos
opuestos. Al ser simples todos los polos, anular el residuo elimina el
polo de la diagonal como divisor. Cada lado exterior aparece una vez.
La fórmula homogénea es de grado cero como forma proyectiva y no tiene
otros divisores polares. Si dos formas racionales proyectivas tuvieran los
mismos residuos y ningún otro polo, su diferencia sería una dos-forma
holomorfa sobre \(\mathbb P^2\), que es necesariamente cero. Así quedan
probadas caracterización y unicidad.
\end{proof}

Se emplea la noción convencional de geometría positiva mediante polos
logarítmicos y residuos canónicos de frontera, desarrollada por
N.~Arkani-Hamed, Y.~Bai y T.~Lam, \emph{Positive Geometries and Canonical
Forms}, JHEP 11 (2017), 039,
\url{https://arxiv.org/abs/1703.04541}, secciones 2, 3, 5.2, 6.1 y 7.2.
La selección de los coeficientes del presente polígono precede a ese
reconocimiento y procede de \(K\).

\subsection{Refinamiento por acarreo sobre la geometría fija}

En un triángulo \(T=[A,B,C]\) de orientación positiva, escribimos
\begin{equation}
 Y(u,r)=(1-u)A+u\{(1-r)B+rC\},\qquad 0<u,r<1.
 \label{eq:polygon-blowdown}
\end{equation}
El vértice \(A\) corresponde a \(u=0\) y el lado opuesto a \(u=1\).
La forma queda
\begin{equation}
 \Omega_T=\frac{du}{u(1-u)}\wedge\frac{dr}{r(1-r)}.
 \label{eq:polygon-product-form}
\end{equation}
Sea \(B_0\geq2\) una base entera declarada, distinta del vértice \(B\).
La coordenada \(r\) admite la actualización por residuo y acarreo
\begin{equation}
 c=\lfloor B_0r\rfloor,\qquad r'=B_0r-c,\qquad
 r=\frac{c+r'}{B_0}.
 \label{eq:polygon-carry}
\end{equation}
Los extremos conservan sus marcas de frontera. El dígito \(c\) permanece
como memoria y permite invertir la actualización. Conviene precisar las
coordenatizaciones de los dos transportes afines. El zoom de la subcelda actúa sobre
\(r\in[c/B_0,(c+1)/B_0]\) mediante \(r'=B_0r-c\).
En cambio, el transporte de la pareja levantada
\(L_c(x,y)=(B_0x-c,B_0y+c)\), con \(x+y=M\), cambia la masa total
a \(B_0M\) y expresa la coordenada
\[
 \frac{B_0x-c}{B_0M}=\frac{x}{M}-\frac{c}{B_0M}.
\]
Son aplicaciones de dominios y normalizaciones diferentes. Ambas conservan
el acarreo en el estado enriquecido y admiten recuperación con sus datos
de escala. La demostración siguiente utiliza el zoom de subcelda y la
partición geométrica que éste identifica.

\begin{theorem}[Invariancia exacta bajo subdivisión con memoria]
\label{thm:polygon-carry}
Sean \(s_j=j/B_0\) y \(W_j=(1-s_j)B+s_jC\).
Los triángulos \(T_j=[A,W_j,W_{j+1}]\), \(0\leq j<B_0\), subdividen
\(T\) y satisfacen
\begin{equation}
 \sum_{j=0}^{B_0-1}\Omega_{T_j}=\Omega_T.
 \label{eq:polygon-carry-additivity}
\end{equation}
La igualdad permanece válida a toda profundidad finita y para cualquier
partición racional ordenada del lado. Subdividir los once triángulos
de \(P_K\) deja \(\Omega_{P_K}\) exactamente invariante.
\end{theorem}
\begin{proof}
Las regiones \(r\in[s_j,s_{j+1}]\) cubren \(T\) con interiores disjuntos.
En las coordenadas comunes,
\[
 \Omega_{T_j}=\frac{du}{u(1-u)}\wedge
 \frac{(s_{j+1}-s_j)\,dr}{(r-s_j)(s_{j+1}-r)}.
\]
La identidad racional
\[
 \frac{s_{j+1}-s_j}{(r-s_j)(s_{j+1}-r)}
 =\frac1{r-s_j}+\frac1{s_{j+1}-r}
\]
telescopa: cada término interior cancela el del subintervalo contiguo.
Quedan \(1/r+1/(1-r)=1/[r(1-r)]\). Esto prueba igualdad de formas
racionales; iterarla proporciona toda profundidad finita.
La prueba sólo utiliza \(s_0=0<s_1<\cdots<s_N=1\), de modo que cubre
particiones no uniformes. El cambio
\(r'=(r-s_j)/(s_{j+1}-s_j)\) devuelve localmente la forma unitaria.
La palabra de acarreos identifica la subcelda y su inversa recupera la
coordenada anterior. Los nuevos polos interiores se cancelan mediante
la misma identidad telescópica.
\end{proof}

La suma anterior reúne la partición geométrica completa. Cuando un lector
de supervivencia selecciona sólo algunas historias, su imagen es la unión
de las subceldas seleccionadas y su forma es la suma correspondiente;
coincide con la forma del polígono completo precisamente cuando esa unión
lo cubre con multiplicidad orientada uno. Por tanto, el criterio de
cobertura para una selección concreta de historias queda expresado por
una condición comprobable sobre su cadena de celdas.

Este refinamiento conserva el polígono fijo. Insertar un nuevo punto
\((t,t^2)\) de la parábola entre dos parámetros consecutivos \(a<t<b\)
produce otra geometría: la altura de la cuerda menos \(t^2\) es
\((t-a)(b-t)>0\), y el nuevo punto queda fuera del polígono previo.
La invariancia demostrada corresponde a la subdivisión rectilínea de
celdas, cuya frontera exterior permanece fija.

\subsection{Cambios de triangulación y reparametrización positiva}

Para cuatro vértices ordenados \(A,B,C,D\), se cumple
\begin{equation}
 \Omega_{ABC}+\Omega_{ACD}=\Omega_{ABD}+\Omega_{BCD}.
 \label{eq:polygon-flip}
\end{equation}
Los dos miembros tienen los mismos residuos exteriores y cancelan su
diagonal interior; su igualdad procede de la unicidad demostrada.
También se comprueba multiplicando por sus denominadores lineales.
Todo cambio entre triangulaciones de un polígono convexo se descompone
en estos cambios de diagonal, por lo que la forma es independiente
de la triangulación. Una identificación con coordenadas cluster debe
conservar adicionalmente la variedad y sus relaciones de cambio.

Si \(D=\operatorname{diag}(d_0,\ldots,d_{12})\), \(d_j>0\), entonces
\begin{equation}
 \mathcal A_{13,1,2}(ZD)=\mathcal A_{13,1,2}(Z).
 \label{eq:polygon-rescale}
\end{equation}
El automorfismo \([c]\mapsto[Dc]\) reparametriza la aplicación cinemática.
En la coordenatización afín, \(c_j\) cambia a
\(c_jd_j/\sum_\ell c_\ell d_\ell\), que recorre nuevamente el símplex.
La imagen y su forma canónica permanecen iguales. El desplazamiento
de una historia parametrizada y el cambio del conjunto geométrico
completo son así operaciones distinguibles.

\subsection{Extensión de rango uno a dimensiones superiores}

Para \(2\leq m\leq12\), definimos
\begin{equation}
 Z_i^{(m)}=(1,t_i,t_i^2,\ldots,t_i^m)^\mathsf T,\qquad
 P_K^{(m)}=\operatorname{conv}
 \{(t_i,t_i^2,\ldots,t_i^m):0\leq i\leq12\}.
 \label{eq:polygon-higher-lift}
\end{equation}
Los menores máximos ordenados son los productos de Vandermonde
\(\prod_{a<b}(t_{i_b}-t_{i_a})>0\). La matriz tiene rango \(m+1\).
La imagen de \(\operatorname{Gr}_{\geq0}(1,13)\) es el polítopo cíclico
\(P_K^{(m)}=\mathcal A_{13,1,m}(Z^{(m)})\), de dimensión \(m\); su fibra
genérica tiene dimensión \(12-m\). Aquí se usa la extensión matemática
de la definición amplituhedral a cualquier \(m\); las aplicaciones
habituales a amplitudes imponen sus propios valores de \(m\).

La triangulación se determina por el orden fijo de etiquetas. Se elige
el menor vértice de un polítopo, se triangulan recursivamente todas sus
facetas que no lo contienen y se toman conos desde ese vértice. La
inducción comienza con los puntos. Una recta desde el vértice inicial
hasta un punto interior corta la frontera opuesta en una faceta; ello
prueba la cobertura. La unicidad del punto de intersección y la
compatibilidad de las restricciones inducidas por el mismo orden prueban
que los interiores son disjuntos y que las intersecciones son caras
comunes. En cada paso disminuye la dimensión, por lo que el procedimiento
termina. Esta es una triangulación por tracción determinada por las
etiquetas.

Para un símplex orientado
\(T=[p_0,\ldots,p_m]\) de esa triangulación, sean
\[
 A_T=(p_1-p_0\ \cdots\ p_m-p_0),\qquad
 (\lambda_1,\ldots,\lambda_m)^\mathsf T=A_T^{-1}(x-p_0),
 \qquad\lambda_0=1-\sum_{j=1}^m\lambda_j.
\]
La orientación de sus vértices se toma de modo que \(\det A_T>0\).
Su forma canónica es
\begin{equation}
 \Omega_T(x)=
 \frac{dx_1\wedge\cdots\wedge dx_m}
 {\det(A_T)\prod_{j=0}^m\lambda_j(x)},\qquad
 \Omega_{P_K^{(m)}}=\sum_{T}\Omega_T.
 \label{eq:polygon-higher-form}
\end{equation}
El cambio baricéntrico prueba la fórmula y que la restricción cinemática
sobre el soporte de \(T\) tiene grado uno. Los residuos son las formas
de las caras con los signos inducidos por la orientación. En una cara
interior, las dos orientaciones son contrarias y el residuo total se
anula. En una cara exterior queda la suma de las formas de la
triangulación de esa cara. La inducción sobre la dimensión demuestra
que dicha suma es su forma canónica. Repetir el residuo a lo largo de
una bandera de caras alcanza los vértices, donde la normalización es
\(\pm1\). Se obtiene así la correspondencia de residuos en todas las
codimensiones para esta colección de rango uno.

La misma demostración cubre cualquier subdivisión simplicial compatible
del polítopo fijo: sus caras interiores se cancelan y las exteriores
conservan los residuos. La inexistencia de formas holomorfas de grado
máximo sobre \(\mathbb P^m\) asegura la unicidad de la forma racional.
En particular, la estabilidad bajo refinamiento, la cobertura y los
residuos quedan cerrados en esta colección de dimensiones \(2,\ldots,12\).
Los grados grassmannianos \(k>1\) conservan sus aplicaciones y sus
obligaciones demostrativas específicas.
